% =============================================================================
% Extrait factuel destiné au rapport présentenciel
% Fichier .tex (fragment) lu par pipeline_rapport_presentenciel.py
%
%   python pipeline_rapport_presentenciel.py --plaidoyer non_coupable \
%       --extrait theories/tex/extrait_faits_rapport_presentenciel.tex
%
% Règles de rédaction de ce fragment :
%   - chaque \section*{...} devient un sous-titre du document final ;
%   - le corps de chaque section est inséré tel quel (LaTeX valide) ;
%   - ne pas utiliser de caractères réservés non échappés (& _ $ % #) ;
%   - ne pas utiliser \section{...} numéroté (le pipeline en ajoute).
% =============================================================================

\section*{Contexte de la cause}

[À COMPLÉTER : identité de l'accusé, juridiction saisie, numéro de dossier,
infraction reprochée, date de l'incident.]

\section*{Faits reprochés}

[À COMPLÉTER : exposé chronologique et factuel des faits allégués par
l'accusation, tel qu'admis ou contesté par la défense.]

\section*{Position de la défense}

[À COMPLÉTER : plaidoyer de non-culpabilité enregistré, moyens de défense
invoqués, contestation des éléments constitutifs de l'infraction.]

\section*{Éléments personnels et sociaux}

[À COMPLÉTER : situation personnelle, familiale, professionnelle, état de
santé, absence ou présence d'antécédents, efforts de réparation le cas
échéant. Ces éléments seront utiles à la détermination de la peine dans
l'éventualité d'un verdict de culpabilité.]

\section*{Pièces et annexes}

[À COMPLÉTER : liste des pièces jointes (ex. : preuves, attestations,
correspondance), avec leur date et leur provenance.]
